\par\noindent\textbf{(1)} 弯曲能为 \(U=\tfrac12EI\int_0^L(y_{xx})^2dx\)，动能 \(K=\tfrac12\mu\int_0^L(y_t)^2dx\)。变分给 \(\mu y_{tt}+EIy_{xxxx}=0\)，是四阶空间偏导方程，边界条件不仅需要位移还可能需要转角或弯矩。
\par\noindent\textbf{(2)} 代入 \(y=e^{i(kx-\omega t)}\) 得 \(\omega^2=(EI/\mu)k^4\)。正频率 \(\omega=\sqrt{EI/\mu}\,k^2\)，相速度 \(v_p=\sqrt{EI/\mu}|k|\)，群速度 \(v_g=2\sqrt{EI/\mu}|k|=2v_p\)。这一理想模型的高频速度无界，说明短波需要Timoshenko修正。
\par\noindent\textbf{(3)} 简支条件是 \(y(0)=y(L)=0\)、\(y_{xx}(0)=y_{xx}(L)=0\)。选择 \(y_n=\sin(n\pi x/L)\) 满足四个边界条件，频率 \(\omega_n=\sqrt{EI/\mu}(n\pi/L)^2\)，所以 \(\omega_n/\omega_1=n^2\)。
\par\noindent\textbf{(4)} 对合法试探形状 \(f(x)\)，\(\omega^2[f]=(EI/\mu)\int(f'')^2dx/\int f^2dx\)。取 \(f=\sin(\pi x/L)\) 直接得到 \(\omega_1^2=(EI/\mu)\pi^4/L^4\)。若改成自由端，则边界为零弯矩与零剪力，不能套用正弦基谱。
